\fig{fig:w2:delay}{同一波形的两个时刻；红点追踪同一相位，虚线标出固定观察位置}{
\begin{tikzpicture}[x=1.25cm,y=.8cm]
\draw[axis](-.2,0)--(4.9,0)node[right]{$x$};
\draw[world,blue,domain=0:4.5,samples=100]plot(\x,{cos(90*\x)});
\fill[red](0,1)circle(2pt)node[above right]{源点$S$在$t_0$的波峰};
\node[left]at(0,0){$0$};
\node at(4.1,1.3){$t=t_0$};
\draw[axis](-.2,-3.4)--(4.9,-3.4)node[right]{$x$};
\draw[world,blue,domain=0:4.5,samples=100]plot(\x,{-3.4+cos(90*(\x-3))});
\fill[red](3,-2.4)circle(2pt)node[above right]{同一波峰到达$P$};
\draw[guide](3,1.45)--(3,-4.6);
\node[below right]at(3,-4.6){$P:\ x$固定};
\draw[dim](0,-4.6)--(3,-4.6)node[midway,below]{$x=u\Delta t$};
\draw[flow](.4,-1.55)--(1.3,-1.55)node[right]{$u$};
\node at(1.1,-2.2){$t=t_0+\Delta t$};
\end{tikzpicture}}

图\ref{fig:w2:delay}比较整条介质在两个时刻的形状。源点$S$与观察点$P$的平衡位置各自固定，红色标记只追踪传播的波峰，不表示同一介质质点从$S$跑到$P$。下图$P$的状态等于上图$S$的状态，因此相同的振动历程到达$P$时整体延迟$x/u$。这一对应关系将在建立行波表达式时使用。
